\chapter{Introduction}
\label{ch:intro}

%-------------------------------------------------------

In Assignment 2 I applied box, Sobel and Gaussian filters, whose kernels
are fixed by formula. A convolutional neural network (CNN) uses the same
sliding-window operation, but it learns the kernel values from labeled
examples \citep{goodfellow2016}. In this assignment I implemented a small CNN in
PyTorch \citep{paszke2019}, trained it to recognize handwritten digits
from the MNIST dataset, followed how it learns batch by batch and epoch
by epoch, and ran a controlled experiment on the learning rate.

%%%%%%%%%%%%%%%%%%%%%%%%%%%%%%%%%%%%%%%%%%%%%%%%%%%%%
\section{Objectives}
\label{sec:objectives}
%%%%%%%%%%%%%%%%%%%%%%%%%%%%%%%%%%%%%%%%%%%%%%%%%%%%%

The assignment has five tasks:

\begin{enumerate}
  \item Implement the CNN in PyTorch and document every layer, its
        input/output dimensions, its trainable parameters and the
        training configuration.
  \item Train the baseline for 10 epochs and plot the training loss,
        accuracy and F1-score batch by batch during the first epoch.
  \item Plot training and validation loss, accuracy and F1-score over
        the 10 epochs and discuss convergence and over/underfitting.
  \item Change one hyperparameter (I chose the learning rate), train the
        alternatives under the same conditions, and pick the best
        configuration using validation data only.
  \item Evaluate the chosen model once on the reserved test set and
        document how to reproduce every result.
\end{enumerate}


%%%%%%%%%%%%%%%%%%%%%%%%%%%%%%%%%%%%%%%%%%%%%%%%%%%%%
\section{Background}
\label{sec:background}
%%%%%%%%%%%%%%%%%%%%%%%%%%%%%%%%%%%%%%%%%%%%%%%%%%%%%

\subsection{Training}

The network outputs 10 scores (logits), one per digit. The cross-entropy
loss turns them into probabilities with softmax and penalizes a low
probability on the correct class:
\begin{equation}
  \mathcal{L} = -\frac{1}{N} \sum_{n=1}^{N} \log
  \frac{e^{z_{n,y_n}}}{\sum_{c=0}^{9} e^{z_{n,c}}}
  \label{eq:ce}
\end{equation}
where $z_{n,c}$ is the score of image $n$ for class $c$ and $y_n$ is its
true label. Training repeats four steps for every mini-batch: forward
pass, loss, back-propagation of the gradients, and a parameter update
\citep{abidl08}. With stochastic gradient descent (SGD) with momentum
$\mu$, the update for parameters $\theta$ with learning rate $\eta$ is
\begin{equation}
  v \leftarrow \mu v + \nabla_\theta \mathcal{L}, \qquad
  \theta \leftarrow \theta - \eta\, v
  \label{eq:sgd}
\end{equation}
The learning rate controls the step size, which is why it is the factor
I study in Task 4.

\subsection{Evaluation Metrics}

Accuracy is the fraction of correctly classified images. For each class
$c$, precision $P_c$ is the fraction of images predicted as $c$ that
really are $c$, and recall $R_c$ is the fraction of images of class $c$
that the model finds. The F1-score is their harmonic mean,
$F1_c = 2 P_c R_c / (P_c + R_c)$. Throughout this report I use
\emph{macro} averages: the metric is computed for each digit separately
and then averaged with equal weight over the set of digits $C$,
\begin{equation}
  \text{macro-F1} = \frac{1}{|C|} \sum_{c \in C} F1_c
  \label{eq:macrof1}
\end{equation}
so a digit the model handles badly lowers the score even if it is rare.
On the full training, validation and test sets, $C$ is all 10 digits.
All three are computed with scikit-learn \citep{pedregosa2011}
(\texttt{average="macro"}).


%%%%%%%%%%%%%%%%%%%%%%%%%%%%%%%%%%%%%%%%%%%%%%%%%%%%%
\section{Report Organization}
\label{sec:organization}
%%%%%%%%%%%%%%%%%%%%%%%%%%%%%%%%%%%%%%%%%%%%%%%%%%%%%

Chapter \ref{ch:setup} describes the dataset, the data split and the
seeding shared by all tasks. Each of the five tasks then has its own
chapter:

\begin{itemize}
  \item Chapter \ref{ch:task1}: Task 1, CNN implementation and model
        documentation;
  \item Chapter \ref{ch:task2}: Task 2, first-epoch batch-level learning
        analysis;
  \item Chapter \ref{ch:task3}: Task 3, ten-epoch training and validation
        analysis;
  \item Chapter \ref{ch:task4}: Task 4, controlled learning-rate
        experiment;
  \item Chapter \ref{ch:task5}: Task 5, independent testing and
        reproducibility.
\end{itemize}

Chapter \ref{ch:conclusion} summarizes the results and their
limitations, and Appendix \ref{app:source} lists the complete source
code.
